% Bagian: first-order-logic
% Bab: syntax-and-semantics
% Subbagian: subformulas

\documentclass[../../../include/open-logic-section]{subfiles}

\begin{document}

\olfileid[id]{fol}{syn}{sbf}

\olsection{\printtoken{P}{subformula}}

\begin{explain}
Sering kali berguna untuk membicarakan !!{formula}s yang ``menyusun'' suatu
!!{formula} tertentu. Kita menyebutnya \emph{!!{subformula}s} dari formula
tersebut. Setiap !!{formula} merupakan !!a{subformula} dari dirinya sendiri;
subformula dari $!A$ selain $!A$ itu sendiri disebut \emph{!!{subformula}
sejati}.
\end{explain}

\begin{defn}[!!^{subformula} langsung]
Jika $!A$ merupakan !!a{formula}, \emph{!!{subformula}s langsung} dari
$!A$ didefinisikan secara induktif sebagai berikut:
\begin{enumerate}
\item !!^{formula}s atomik tidak memiliki !!{subformula} langsung.

\tagitem{prvNot}{\indcase{!A}{\lnot !B}{Satu-satunya !!{subformula}
    langsung dari $\indfrm$ ialah~$!B$.}}{}

\item \indcase{!A}{(!B \ast !C)}{!!^{subformula}s langsung dari
  $\indfrm$ ialah $!B$ dan $!C$ ($\ast$ merupakan salah satu penghubung
  dua-tempat).}

\tagitem{prvAll}{\indcase{!A}{\lforall[x][!B]}{Satu-satunya
    !!{subformula} langsung dari $\indfrm$ ialah~$!B$.}}{}

\tagitem{prvEx}{\indcase{!A}{\lexists[x][!B]}{Satu-satunya
    !!{subformula} langsung dari $\indfrm$ ialah~$!B$.}}{}
\end{enumerate}
\end{defn}

\begin{defn}[!!^{subformula} sejati]
Jika $!A$ merupakan !!a{formula}, \emph{!!{subformula}s sejati} dari
$!A$ didefinisikan secara rekursif sebagai berikut:
\begin{enumerate}
\item !!^{formula}s atomik tidak memiliki !!{subformula} sejati.

\tagitem{prvNot}{\indcase{!A}{\lnot !B}{!!^{subformula}s sejati dari
    $\indfrm$ ialah~$!B$ bersama semua !!{subformula}s sejati
    dari~$!B$.}}{}

\item \indcase{!A}{(!B \ast !C)}{!!^{subformula}s sejati dari
  $\indfrm$ ialah $!B$ dan $!C$, bersama semua !!{subformula}s sejati
  dari $!B$ dan semua !!{subformula}s sejati dari~$!C$.}

\tagitem{prvAll}{\indcase{!A}{\lforall[x][!B]}{!!^{subformula}s sejati
    dari $\indfrm$ ialah~$!B$ bersama semua !!{subformula}s sejati
    dari~$!B$.}}{}

\tagitem{prvEx}{\indcase{!A}{\lexists[x][!B]}{!!^{subformula}s sejati
    dari $\indfrm$ ialah~$!B$ bersama semua !!{subformula}s sejati
    dari~$!B$.}}{}
\end{enumerate}
\end{defn}

\begin{defn}[!!^{subformula}]
!!^{subformula}s dari $!A$ ialah $!A$ itu sendiri bersama semua
!!{subformula}s sejatinya.
\end{defn}

\begin{explain}
Perhatikan perbedaan halus antara cara kita mendefinisikan !!{subformula}s
langsung dan !!{subformula}s sejati. Dalam kasus pertama, kita secara
langsung mendefinisikan !!{subformula}s langsung suatu formula~$!A$ untuk
setiap bentuk yang mungkin dari~$!A$. Definisi tersebut merupakan definisi
eksplisit menurut kasus, dan kasus-kasusnya mencerminkan definisi induktif
himpunan !!{formula}s. Dalam kasus kedua, kita juga mencerminkan cara
himpunan semua !!{formula}s didefinisikan, tetapi dalam setiap kasus kita
turut menyertakan !!{subformula}s sejati dari !!{formula}s yang lebih kecil,
yaitu $!B$ dan $!C$, selain !!{formula}s tersebut sendiri. Hal ini membuat
definisi itu \emph{rekursif}. Secara umum, definisi suatu fungsi pada
himpunan yang didefinisikan secara induktif (dalam kasus kita,
!!{formula}s) bersifat rekursif jika kasus-kasus dalam definisi fungsi itu
menggunakan fungsi tersebut sendiri. Namun, agar terdefinisi dengan baik,
kita harus memastikan bahwa kita hanya menggunakan nilai fungsi pada
argumen yang muncul ``sebelum'' argumen yang sedang kita definisikan---dalam
kasus kita, ketika mendefinisikan ``!!{subformula} sejati'' untuk $(!B
\ast !C)$, kita hanya menggunakan !!{subformula}s sejati dari
!!{formula}s yang ``lebih dahulu,'' yaitu $!B$ dan $!C$.
\end{explain}

\begin{prop}
\ollabel{prop:subfrm-trans}
Andaikan $!B$ merupakan subformula dari $!A$ dan $!C$ merupakan subformula
dari $!B$. Maka $!C$ merupakan subformula dari $!A$. Dengan kata lain,
relasi subformula bersifat transitif.
\end{prop}

\begin{prob}
Buktikan \olref[fol][syn][sbf]{prop:subfrm-trans}.
\end{prob}

\begin{prop}
\ollabel{prop:count-subfrms}
Andaikan $!A$ merupakan formula dengan $n$ penghubung dan kuantor. Maka
$!A$ memiliki paling banyak $2n+1$ subformula.
\end{prop}

\begin{prob}
Buktikan \olref[fol][syn][sbf]{prop:count-subfrms}.
\end{prob}

\end{document}

